\subsection{Inductive Voltage Divider}

This block is constructed from two inductors in series and is used
to attenuate a voltage signal.

\subsubsection{Schematic}
\begin{center}
  \begin{tikzpicture}[transform shape]
	  % Paths, nodes and wires:
	  \node[ground, rotate=-90] at (5.5, 5.25){};
	  \node[ground, rotate=-90] at (5.5, 4.75){};
	  \draw (5.5, 6.75) to[sinusoidal voltage source, l={$V_{cc}$}] (5.5, 5.25);
	  \draw (5.5, 4.75) to[sinusoidal voltage source, l={$V_{ee}$}] (5.5, 3.25);
	  \draw (7, 5) -- (7.75, 5);
	  \draw (7, 7) -| (5.5, 7) -| (5.5, 6.75);
	  \node[shape=rectangle, minimum width=0.715cm, minimum height=0.465cm](N1) at (8.125, 5){} node[anchor=center] at (N1.text){$V_{out}$};
	  \draw (7, 5.25) -| (7, 5);
	  \draw (7, 4.75) -| (7, 5);
	  \draw (5.5, 3.25) -| (5.5, 3) -- (7, 3);
	  \draw (7, 7) to[american inductor, l={$L_1$}] (7, 5.25);
	  \draw (7, 4.75) to[american inductor, l={$L_1$}] (7, 3);
  \end{tikzpicture}
\end{center}

\subsubsection{Transfer Function}
Applying KCL at $V_{out}$ we get:
\[
  \frac{V_{out} - V_{cc}}{sL_1} + \frac{V_{out} - V_{ee}}{sL_2} = 0
\]

\paragraph{Superposition:}
\begin{enumerate}
  \item Due to \(V_{cc}\) alone (\(V_{ee} = 0V\)):
    \begin{align*}
      \frac{V_{out} - V_{cc}}{sL_1} + \frac{V_{out} - 0}{sL_2} &= 0 \\
      \frac{V_{out} - V_{cc}}{sL_1} + \frac{V_{out}}{sL_2} &= 0 \\
      \frac{V_{out}}{sL_1} - \frac{V_{cc}}{sL_1} + \frac{V_{out}}{sL_2} &= 0 \\
      V_{out}\left(\frac{1}{sL_1} + \frac{1}{sL_2}\right) - \frac{V_{cc}}{sL_1} &= 0 \\
      V_{out}\left(\frac{L_2 + L_1}{sL_1 L_2}\right) - \frac{V_{cc}}{sL_1} &= 0 \\
      V_{out}\left(\frac{L_2 + L_1}{sL_1 L_2}\right) &= \frac{V_{cc}}{sL_1} \\
      V_{out}(L_2 + L_1) &= V_{cc} L_2 \\
      V_{out} &= V_{cc} \frac{L_2}{L_1 + L_2} \\
      \boxed{H_1(s) = \frac{N_1(s)}{D_1(s)} = \frac{V_{out_1}}{V_{cc}} = \frac{L_2}{L_1 + L_2}}
    \end{align*}
  \item Due to \(V_{ee}\) alone (\(V_{cc} = 0V\)):
    \begin{align*}
      \frac{V_{out} - 0}{sL_1} + \frac{V_{out} - V_{ee}}{sL_2} &= 0 \\
      \frac{V_{out}}{sL_1} + \frac{V_{out} - V_{ee}}{sL_2} &= 0 \\
      \frac{V_{out}}{sL_1} + \frac{V_{out}}{sL_2} - \frac{V_{ee}}{sL_2} &= 0 \\
      V_{out}\left(\frac{1}{sL_1} + \frac{1}{sL_2}\right) - \frac{V_{ee}}{sL_2} &= 0 \\
      V_{out}\left(\frac{L_2 + L_1}{sL_1 L_2}\right) - \frac{V_{ee}}{sL_2} &= 0 \\
      V_{out}\left(\frac{L_2 + L_1}{sL_1 L_2}\right) &= \frac{V_{ee}}{sL_2} \\
      V_{out}(L_2 + L_1) &= V_{ee} L_1 \\
      V_{out} &= V_{ee} \frac{L_1}{L_1 + L_2} \\
      \boxed{H_2(s) = \frac{N_2(s)}{D_2(s)} = \frac{V_{out_2}}{V_{ee}} = \frac{L_1}{L_1 + L_2}}
    \end{align*}
\end{enumerate}
\paragraph{Total Output:}
\begin{align*}
  V_{out} &= V_{out_1} + V_{out_2} \\
  V_{out} &= V_{cc} \frac{L_2}{L_1 + L_2} + V_{ee} \frac{L_1}{L_1 + L_2}
\end{align*}


\begin{enumerate}
  \item \textbf{Inductor Ratio (\(\frac{L_1}{L_2}\)):}
    From KCL at \(V_{out}\):
    \begin{align*}
      \frac{V_{out} - V_{cc}}{sL_1} - \frac{V_{ee} - V_{out}}{sL_2} &= 0 \\
      \frac{V_{out} - V_{cc}}{sL_1} &= \frac{V_{ee} - V_{out}}{sL_2} \\
      \frac{sL_2}{sL_1} &= \frac{V_{ee} - V_{out}}{V_{out} - V_{cc}} \\
      \boxed{\frac{L_2}{L_1} = \frac{V_{ee} - V_{out}}{V_{out} - V_{cc}}}
    \end{align*}
  \item \textbf{Impedance's:}
    \paragraph{At Input:}
      \begin{align*}
        Z_{in} &= sL_1 + sL_2 \\
        \boxed{Z_{in} = s(L_1 + L_2)}
      \end{align*}
    \paragraph{At Output:}
      \begin{align*}
        Z_{out} &= \frac{1}{\frac{1}{sL_1} + \frac{1}{sL_2}} \\
        Z_{out} &= \frac{1}{\frac{L_2 + L_1}{sL_1 L_2}} \\
        \boxed{Z_{out} = \frac{sL_1 L_2}{L_1 + L_2}}
      \end{align*}
  \item \textbf{Loading Effect:} \\
    Similar to resistive dividers.
\end{enumerate}

\subsubsection{Question}
Design an inductive voltage divider that attenuates a 10V signal to 7V.
Establish the ratio of the inductors and calculate the input and output
impedance. If one of the inductors is 18 nH, calculate the other inductor's
value.
\paragraph{Solution:}
\begin{enumerate}
  \item Calculate the inductor ratio:
    \begin{align*}
      \frac{L_2}{L_1} &= \frac{V_{ee} - V_{out}}{V_{out} - V_{cc}} \\
      \frac{L_2}{L_1} &= \frac{0V - 7V}{7V - 10V} \\
      \frac{L_2}{L_1} &= \frac{-7V}{-3V} \\
      \frac{L_2}{L_1} &= \frac{7}{3} \\
      \boxed{\frac{L_2}{L_1} = 2.33}
    \end{align*}
  \item Calculate the input impedance:
    \begin{align*}
      Z_{in} &= s(L_1 + L_2) \\
      Z_{in} &= s(L_1 + 2.33L_1) \hspace{1cm} \text{[substituting \(L_2 = 2.33L_1\)]} \\
      Z_{in} &= sL_1(1 + 2.33) \\
      Z_{in} &= s(3.33L_1)
    \end{align*}
  \item Calculate the output impedance:
    \begin{align*}
      Z_{out} &= \frac{sL_1 L_2}{L_1+L_2} \\
      Z_{out} &= \frac{sL_1(2.33L_1)}{L_1 + 2.33L_1} \hspace{1cm} \text{[substituting \(L_2 = 2.33L_1\)]} \\
      Z_{out} &= \frac{s(2.33L_1^2)}{L_1(1+2.33)} \\
      Z_{out} &= \frac{s(2.33L_1^2)}{3.33L_1} \\
      Z_{out} &= s(0.699L_1)
    \end{align*}
  \item Calculate the other inductor's value:
    \begin{align*}
      L_2 &= 2.33L_1 \\
      L_2 &= 2.33(18nH) \\
      L_2 &\approx 41.94nH
    \end{align*}
\end{enumerate}


